% ------------------------------
% Header LaTeX para Quarto (PDF)
% Proyecto Docente — estilo sobrio/profesional
% ------------------------------
\typeout{=== HEADER.TEX CARGADO ===}

% Colores (necesario para \textcolor y \definecolor)
\usepackage{xcolor}
\definecolor{UAMBlue}{HTML}{002060} % azul institucional (#002060)

% Microtipografía (mejora “acabado” del texto)
\usepackage{microtype}

% Evita que LaTeX estire verticalmente las páginas
\raggedbottom

% --- Cabeceras/pies con KOMA (recomendado en Quarto PDF) ---
\usepackage[automark]{scrlayer-scrpage}
\clearpairofpagestyles

% Estilo: pequeño, cursiva, gris
\newcommand{\HeaderStyle}{\footnotesize\itshape\color{gray}}

% Cabeceras pares/impares
\lehead{\HeaderStyle Modesto Redrejo Rodríguez} % páginas pares
\rohead{\HeaderStyle Proyecto Docente Plaza UAM2026-A1055TU}          % páginas impares

% Línea bajo cabecera (opcional)
\setheadsepline{0.3pt}

% Número de página en pie (también en páginas 'plain')
\cfoot*{\footnotesize\color{gray}\pagemark}

%spanish
\usepackage[spanish]{babel}
\renewcommand{\tablename}{Tabla}
\renewcommand{\figurename}{Figura}

% ACTIVAR el estilo con cabecera
\pagestyle{scrheadings}

% IMPORTANTÍSIMO: los inicios de capítulo suelen usar 'plain' (sin cabecera).
% Forzamos que también usen scrheadings:
\renewcommand*{\chapterpagestyle}{scrheadings}

% Seguro extra: por si el template cambia el pagestyle más tarde
\AtBeginDocument{\pagestyle{scrheadings}}

% Tablas y figuras: mejor aspecto tipográfico
\usepackage{booktabs}
\usepackage{caption}
\captionsetup{
  font=small,
  labelfont=bf
}

% Evitar viudas/huérfanas (detalle de “calidad editorial”)
\widowpenalty=10000
\clubpenalty=10000


%formato tablitas
\usepackage{array}
\setlength{\arrayrulewidth}{0.3pt}
\setlength{\doublerulesep}{0pt}
%\usepackage{booktabs}
%\setlength{\lightrulewidth}{0.3pt}
%\renewcommand{\arraystretch}{1.2}


\AtBeginDocument{
  \renewcommand{\tablename}{Tabla}
  \renewcommand{\figurename}{Figura}
}

\usepackage[table]{xcolor}

% Define gris muy suave
\definecolor{lightgray}{gray}{0.95}

\usepackage{enumitem}
\setlist[itemize]{leftmargin=2em, noitemsep, topsep=0pt, partopsep=0pt, parsep=0pt}

\usepackage{longtable}
\usepackage{multirow}
\usepackage{graphicx}
\usepackage{float}
\usepackage{array} 
\usepackage[table]{xcolor}
\usepackage{hyperref}
\usepackage{fontawesome5}
